\section{An encoding upper bound in continuous dimension}\label{sec:upper}
We now consider rational quadratic input of total binary length $N\ge2$:
\begin{equation}\label{eq:quadraticmodel}
 v=\min\{f(x,z):(x,z)\in X,\ Ax+Bz-b=0\},\qquad
 x\in\R^n,\ z\in\Z^p,
\end{equation}
where $n\ge1$ and the residual $r(x,z)=Ax+Bz-b$ has $m\ge1$
coordinates. The native set is
specified by finite rational coordinate bounds and inequalities
$g_i(x,z)\le0$, all of total degree at most two; $f$ has the same degree
bound. Let $s$ count all native inequalities, including the boxes. Affine
equalities may be represented by paired inequalities. The input lists
expanded polynomials, with coefficients and variable indices encoded
explicitly; in particular, $s,m,n\le N$.

For every integer assignment, assume that $f(\cdot,z)$ and each
$g_i(\cdot,z)$ are convex on $\R^n$. Write
$X_z=\{x:(x,z)\in X\}$. A slice is \emph{equality-feasible} if
$X_z\cap\{x:Ax+Bz=b\}\ne\varnothing$.
Suppose the original feasible set is nonempty and, on every such slice,
this intersection contains a point strictly satisfying every native
inequality nonlinear in $x$; native affine inequalities may hold weakly.
This is the \emph{refined Slater condition}. No margin is given,
and no Slater condition is imposed on equality-infeasible slices.
Empty native slices may be discarded. The finite boxes leave only finitely
many integer assignments.

Write $\operatorname{bit}(a)=1+\lfloor\log_2 a\rfloor$ for a positive
integer $a$. All constants implicit below are absolute and independent of
the integer dimension and the number of slices.

\begin{theorem}\label{thm:generalupper}
Under these assumptions there is an integer $\rho\ge1$ with
\begin{equation}\label{eq:generalbits}
 \operatorname{bit}(\rho)\le (N+1)2^{O(n)}
\end{equation}
such that $(\rho,0)$ is minimizer-set exact for either
$\norm{\cdot}_\infty$ or $\norm{\cdot}_1$.
\end{theorem}
We use the following consequence of effective bounds for real algebraic
sets. A \emph{basic closed set} is a conjunction of polynomial equalities
and weak inequalities.

\begin{lemma}[Effective radii]\label{lem:radii}
Let at most $S$ integer polynomials in $d$ variables have degree at most
two and coefficient bitsize at most $\tau\ge1$. There is a radius $R\ge1$
with
\begin{equation}\label{eq:radius}
 \log_2 R\le(\tau+\log(S+1)+1)2^{O(d)}
\end{equation}
that meets every nonempty basic closed set defined by these polynomials
and contains every such set that is bounded. The same statements hold
for a finite union of such sets.
\end{lemma}
\begin{proof}
Apply the component-meeting and bounded-component radii of
\citet[Theorems~3--4]{BasuRoy2010} to weak sign conditions and take the
larger radius. Their explicit formulas give \eqref{eq:radius} for fixed
degree. A conjunction that leaves some polynomial signs unspecified is
a finite union of compatible weak sign conditions. For a union, select
one nonempty member to obtain a point; if the union is bounded, every
member and every component is bounded. The source proof's factorial-bit
term, retained when majorizing its meeting radius, is also absorbed by
$2^{O(d)}$.
\end{proof}
Only closed conjunctions and their finite unions are used here; the lemma
makes no claim about arbitrary strict sign conditions.

\subsection{Uniform coefficient bounds}
The explicit input gives coordinate bounds of magnitude $2^{O(N)}$ and
$O(N)$ encoded monomials. Multiplying by the product of all input
denominators clears coefficients with $O(N)$ bits. Substituting any boxed
integer assignment into a degree-two monomial adds only $O(N)$ bits;
summing monomials adds $O(\log N)$ more. Differentiation and introducing
multiplier variables preserve this uniform $O(N)$ coefficient-height
bound. Thus every degree-two polynomial system below can be cleared with
\begin{equation}\label{eq:height}
 \tau=O(N).
\end{equation}
We clear the polynomial equations of certificate systems, not the
residual in the penalty; the multiplier coordinates keep their original
meaning. A box bound $|f|\le F_0$, with an integer $F_0\ge1$ and
$\operatorname{bit}(F_0)=O(N)$, is equally immediate. For example, if
$H\ge1$ bounds all coordinate magnitudes and $C$ is the sum of the
absolute objective coefficients, take $F_0=1+\lceil CH^2\rceil$.

\subsection{Feasible slices: a sparse multiplier certificate}
Fix an equality-feasible assignment and suppress $z$. Refined Slater
ensures a convex KKT certificate $(x^*,\mu,\lambda)$, with
$\mu\ge0$. To apply the usual relative-interior form of Slater, one may
blend the strictly feasible point with a point in the relative interior
of the polyhedron defined by the affine constraints. For a sufficiently
small blend, all nonlinear inequalities remain strict. The resulting
convex optimality conditions express the normal of that polyhedron by
its active affine rows, giving the stated KKT certificate even when
these rows are dependent.

Let $e=\operatorname{rank} A$ and select independent rows $J$ spanning
its row space. Project stationarity onto the quotient by this row space.
The projected vector $-\nabla f(x^*)$ is in the cone generated by active
native gradients, in dimension $n-e$. Conic Carath\'eodory therefore
selects at most $n-e$ gradients, indexed by $I$. The remaining vector
is in the row space of $A$. Hence $|I|+|J|\le n$, and the following
system has a solution:
\begin{equation}\label{eq:sparsekkt}
 \begin{aligned}
 &g_i(x)\le0\quad(1\le i\le s),\qquad Ax+Bz-b=0,\\
 &\mu_i\ge0,\quad g_i(x)=0\quad(i\in I),\\
 &\nabla f(x)+\sum_{i\in I}\mu_i\nabla g_i(x)
                  +\sum_{j\in J}\lambda_j A_j^T=0.
 \end{aligned}
\end{equation}
This is a basic closed set in at most $2n$ variables, defined by at most
$s+m+2n$ degree-two polynomials: the active equations reuse the native
polynomials $g_i$. Replacing complementarity products by these equations
keeps degree two. The selected original
certificate proves nonemptiness; every point in the new system remains
a KKT certificate.

By Lemma~\ref{lem:radii}, some certificate has all coordinates bounded
by $R_M$, where
\begin{equation}\label{eq:RM}
 \log_2 R_M\le (\tau+\log(s+m+2n+1)+1)2^{O(n)}.
\end{equation}
Set omitted linking multipliers to zero. The resulting multiplier has
$\norm{\lambda_z}_1\le nR_M$. Convexity, stationarity and
complementarity imply, throughout the native slice,
\begin{equation}\label{eq:slicemultiplier}
 f(x,z)+\lambda_z^T(Ax+Bz-b)\ge v_z.
\end{equation}
Indeed, the full convex Lagrangian is minimized at the KKT point and
has value $v_z$ there; deleting its nonpositive native-inequality terms
on $X_z$ only increases it. We bound one suitable certificate, not all
optimal multipliers, whose set can be unbounded.

\subsection{Infeasible slices: a reciprocal graph}
For a nonempty equality-infeasible native slice, compactness gives
$\delta_z=\min_{x\in X_z}\norm{Ax+Bz-b}_\infty>0$.
Write $r_j=r_j(x,z)$ and consider
\begin{equation}\label{eq:reciprocal}
 T_z=\left\{(x,t):\begin{array}{l}
 x\in X_z,\quad t\ge0,\quad -1\le tr_j\le1\ (1\le j\le m),\\
 tr_j=1\ \text{or}\ tr_j=-1\ \text{for at least one }j
 \end{array}\right\}.
\end{equation}
This is exactly the graph of $t=1/\norm{r(x,z)}_\infty$. It is
nonempty and compact, since the denominator is bounded away from zero
on $X_z$. Selecting a coordinate and a sign in the last line expresses
it as a union of $2m$ basic closed sets. The equality uses one of the
polynomials already present in $-1\le tr_j\le1$, so only
$s+2m+1$ degree-two polynomials in $n+1$ variables are needed.
The bounded-set part of Lemma~\ref{lem:radii} gives
\begin{equation}\label{eq:RD}
 \delta_z^{-1}\le R_D,\qquad
 \log_2 R_D\le(\tau+\log(s+2m+2)+1)2^{O(n)}.
\end{equation}
No KKT condition or rationality of $\delta_z$ is used on these slices.

\begin{proof}[Proof of Theorem~\ref{thm:generalupper}]
Choose an integer $E\ge0$ with $2^E\ge\max\{R_M,R_D\}$ satisfying
the preceding bounds, and set
\begin{equation}\label{eq:sufficientrho}
 \rho=(n+2F_0+1)2^E.
\end{equation}
This is strictly larger than both $nR_M$ and $2F_0R_D$.
Lemma~\ref{lem:slices}, with objective bounds $-F_0,F_0$, proves
minimizer-set exactness for the infinity norm. In detail, on a feasible
slice a point of nonzero residual has penalized value at least
$v_z+(\rho-nR_M)\norm r_\infty>v$, while on an infeasible slice it
has value at least $-F_0+\rho/R_D>F_0\ge v$.
The same coefficient works for the one-norm because
$\norm r_1\ge\norm r_\infty$ and both vanish on the original feasible
set. Equations \eqref{eq:height}--\eqref{eq:RD} and
$\operatorname{bit}(F_0)=O(N)$ give \eqref{eq:generalbits}. All bounds are
uniform over boxed integer assignments, so no factor counting those
assignments appears.
\end{proof}
Together with Theorem~\ref{thm:lower}, this gives exponential upper
and lower dependence of penalty bit length on continuous dimension.
The constants in the exponents and the dependence on the full sparse
input length are not asserted to be sharp.
